\subsection{Local moment relaxations}\label{app:moments}

This subsection gives the details behind Section~\ref{sec:limits-local}.
Throughout, $F(x)=\frac12x^\T Hx+b^\T x+F(0)$ is a quadratic on a continuous box
$X$ with $H\succeq-\nu I$ for some $\nu\ge0$. An error bound in which $L$ is
replaced by $\nu$ holds when all coordinates are conditioned on one global
cell.

\begin{lemma}[Globally conditioned moments]\label{lim:lem:mixture}
Let $\bar x\in C=\prod_i[\ell_i',u_i']\subseteq X$, and let $Y$ be a symmetric
matrix with $Y-\bar x\bar x^\T\succeq0$ and
$Y_{ii}\le(\ell_i'+u_i')\bar x_i-\ell_i'u_i'$ for all $i$. Then
\[
\tfrac12\langle H,Y\rangle+b^\T\bar x+F(0)\ge\OPT-\frac\nu8\sum_i(u_i'-\ell_i')^2 .
\]
The same bound, averaged with the corresponding weights, holds for any
convex combination of such pairs $(\bar x,Y)$ indexed by cells $C$.
\end{lemma}

\begin{proof}
Put $\Sigma=Y-\bar x\bar x^\T\succeq0$. The left side equals
$F(\bar x)+\frac12\langle H,\Sigma\rangle\ge F(\bar x)-\frac\nu2\operatorname{tr}\Sigma$.
Also $\Sigma_{ii}=Y_{ii}-\bar x_i^2\le(u_i'-\bar x_i)(\bar x_i-\ell_i')
\le(u_i'-\ell_i')^2/4$, and $F(\bar x)\ge\OPT$. Averaging preserves the
inequality.
\end{proof}

A global cell index can take $K^n$ values when every coordinate is split into
$K$ intervals. Sparse relaxations keep exact information only on bags and
match moments across separators.

\begin{definition}[Local moment relaxation of order $k$]\label{lim:def:moments}
Fix a tree decomposition and a finite partition of every $X_i$ into intervals
(cells). A feasible point $\mu$ consists of probability measures $\mu_t$ on
$X_{B_t}=\prod_{i\in B_t}X_i$, one per bag, such that for every tree edge
$\{t,u\}$, every product $C$ of cells of the coordinates in $S_{tu}$, and
every monomial $x_{S_{tu}}^\alpha$ with $\abs\alpha\le k$,
\[
\int_{\{x_{S_{tu}}\in C\}}x_{S_{tu}}^\alpha\,d\mu_t
=\int_{\{x_{S_{tu}}\in C\}}x_{S_{tu}}^\alpha\,d\mu_u .
\]
Its value is $\sum_t\int q_t\,d\mu_t$, where $q_t$ is the sum of the factors
assigned to $t$. Point masses at points of $X$ are feasible, so the optimal
value of the relaxation is at most $\OPT$. The \emph{occupied width}
$\omega_i(\mu)$ is the largest length of a cell of coordinate $i$ that carries
positive mass under some $\mu_t$ with $i\in B_t$.
\end{definition}

\begin{proposition}[Every finite order fails]\label{lim:prop:moments}
For each integer $r\ge1$ and rational $0<h\le1/(4r)$ there is a continuous box
QP $F_{r,h}$ in $2r-1$ state variables with range $[-1,1]$ and $2r-1$ control
variables with range $[0,1]$, with a tree decomposition of maximum bag size
three, a unique minimizer and $\OPT=1/4$, such that
$\nabla^2F_{r,h}\succeq-2I$ and $g_r=1/(8r(1+8(2r-1)^2))$ is a valid growth
constant, both independent of $h$, with the following property. For every
choice of cells in which $[0,rh]$ lies in one cell of length at most $2rh$ for
each state variable, and $0$ and $1$ lie in cells of length at most $h'$ for
each control variable, the relaxation of order $2r-1$ has a feasible point
$\mu$ with value $\frac14-\frac{2r+1}{16r}$ and
\[
\sum_i\omega_i(\mu)^2\le(4r-2)\max\{2rh,h'\}^2 .
\]
Consequently there is no function $\phi$ such that every feasible point $\mu$
of the relaxation of order $k$ has value at least
$\OPT-\phi(p,\bar\nu/g,k)\,\bar\nu\sum_i\omega_i(\mu)^2$, for all instances,
all $\bar\nu\ge0$ with $\nabla^2F\succeq-\bar\nu I$, all growth constants $g$,
all orders $k$ and all cells.
\end{proposition}

\begin{proof}
\emph{Construction.} Use state variables $s$, $y_1,\ldots,y_{r-1}$,
$z_1,\ldots,z_{r-1}$ in $[-1,1]$ and control variables $u_1,\ldots,u_r$,
$v_1,\ldots,v_{r-1}$ in $[0,1]$. Put $y_0=z_0=0$, $y_r=s$, and
\[
\rho_i=\frac{y_i-y_{i-1}}h-u_i\ (1\le i\le r),\quad
\rho'_j=\frac{z_j-z_{j-1}}h-v_j\ (1\le j<r),\quad
\rho'_r=\frac{s-z_{r-1}}h-\frac12,
\]
\[
F_{r,h}=2r\Bigl(\sum_{i=1}^r\rho_i^2+\sum_{j=1}^r{\rho'_j}^2\Bigr)
+\sum_{i=1}^ru_i(1-u_i)+\sum_{j=1}^{r-1}v_j(1-v_j)+\lambda(\bar u+\bar v),
\]
where $\bar u=\sum_iu_i$, $\bar v=\sum_jv_j$ and $\lambda=1/(8r)$. The bags
$B^{\mathrm L}_i=\{y_{i-1},y_i,u_i\}$ ($i\le r$),
$B^{\mathrm R}_j=\{z_{j-1},z_j,v_j\}$ ($j<r$) and
$B^{\mathrm R}_r=\{z_{r-1},s\}$, with constants omitted, are arranged in the
path $B^{\mathrm L}_1,\ldots,B^{\mathrm L}_r,B^{\mathrm R}_r,
B^{\mathrm R}_{r-1},\ldots,B^{\mathrm R}_1$. The separator between
$B^{\mathrm L}_r$ and $B^{\mathrm R}_r$ is $\{s\}$, and all bags have at most
three elements.

\emph{Reduction.} List the rescaled path nodes $\zeta_0=0$, $\zeta_k=y_k/h$
($1\le k<r$), $\zeta_r=s/h$, $\zeta_{r+l}=z_{r-l}/h$ ($1\le l<r$),
$\zeta_{2r}=0$. Up to sign, the $2r$ residuals are $\zeta_k-\zeta_{k-1}-c_k$
with $c=(u_1,\ldots,u_r,-\frac12,-v_{r-1},\ldots,-v_1)$; note
$\bar c:=\sum_kc_k=\bar u-\bar v-\frac12$. Let $\zeta^0(u,v)$ be the node values
with every residual equal to $-\bar c/(2r)$, and $\delta=\zeta-\zeta^0$, with
$\delta_0=\delta_{2r}=0$. Since the residuals sum to $-\bar c$, expanding the
squares gives
\begin{equation}\label{lim:eq:momsplit}
F_{r,h}=\Psi(u,v)+2r\sum_{k=1}^{2r}(\delta_k-\delta_{k-1})^2,\quad
\Psi=\bar c^2+\sum_iu_i(1-u_i)+\sum_jv_j(1-v_j)+\lambda(\bar u+\bar v).
\end{equation}
Since $\abs{\bar c}<r$ and $\abs{\sum_{l\le k}c_l}\le r$, every node of
$\zeta^0$ lies in $[-2r,2r]$. Hence the corresponding state variables
$h\zeta^0$ lie in $[-\tfrac12,\tfrac12]$ and are feasible.

\emph{The reduced function.} With $e_2$ the second elementary symmetric
polynomial,
\begin{equation}\label{lim:eq:pi}
\Psi-\frac14=2\Upsilon+\lambda(\bar u+\bar v),\qquad
\Upsilon=e_2(u)+e_2(v)-\bar u\bar v+\bar v .
\end{equation}
This follows from $\bar c^2=(\bar u-\bar v)^2-(\bar u-\bar v)+\frac14$,
$\sum_iu_i(1-u_i)=\bar u-\sum_iu_i^2$, and $\bar u^2-\sum_iu_i^2=2e_2(u)$. The
polynomial $\Upsilon$ is affine in each variable, so its minimum over the box
is attained at a vertex. At a vertex with $\bar u=a$ and $\bar v=b$,
$\Upsilon=\frac12[(a-b)^2-(a-b)]=\frac12(a-b)(a-b-1)\ge0$, because $a-b$ is an
integer. Therefore $\Psi\ge\frac14+\lambda(\bar u+\bar v)$, and
\eqref{lim:eq:momsplit} shows that the unique minimizer is $u=v=0$,
$\zeta=\zeta^0(0,0)$, with $\OPT=\frac14$.

\emph{Growth and curvature.} By \eqref{lim:eq:momsplit}, \eqref{lim:eq:pi},
$\bar u+\bar v\ge\norm{(u,v)}^2$, and the path inequality \eqref{lim:eq:path}
with $2r$ in place of $m$,
\[
F_{r,h}-\tfrac14\ge\lambda\norm{(u,v)}^2+\frac2{2r-1}\norm\delta^2 .
\]
Each node of $\zeta^0$ is a partial sum of the $c_l$ minus $k/(2r)$ times
$\bar c$, so each of the $2r-1$ free nodes moves by at most
$2\norm{(u,v)}_1\le2\sqrt{2r-1}\,\norm{(u,v)}$ when the controls move from
$0$ to $(u,v)$. Hence
$\norm{\zeta^0(u,v)-\zeta^0(0,0)}^2\le4(2r-1)^2\norm{(u,v)}^2$. The state
variables are $h$ times the nodes and $h\le1$, so the squared distance to the
minimizer is at most $2\norm\delta^2+(1+8(2r-1)^2)\norm{(u,v)}^2$. This gives
the growth constant $g_r=\lambda/(1+8(2r-1)^2)\le1/(2r-1)$. The Hessian is $4r$
times a sum of rank-one Gram terms minus $2$ on the controls, so
$\nabla^2F_{r,h}\succeq-2I$.

\emph{The relaxation witness.} On the left subtree take the law under which
$\bar u=i$ with probability $\pi_i=\binom{2r}{2i}2^{1-2r}$ ($0\le i\le r$),
$u=(1^i,0^{r-i})$, $y_k=h\sum_{l\le k}u_l$ and $s=hi$. On the right subtree
take $\bar v=j$ with probability $\pi'_j=\binom{2r}{2j+1}2^{1-2r}$ ($0\le j<r$),
$v=(1^j,0^{r-1-j})$, $z_k=h\sum_{l\le k}v_l$ and $s=h(j+\frac12)$. Both are
probability laws, all residuals vanish on their supports, and all controls are
binary. Let $\mu_t$ be the corresponding bag marginals. At every separator
other than $\{s\}$ the two bag marginals come from one law, so they agree on
every set and in all moments. At the separator $\{s\}$ the laws of $2s/h$ are
those of a $\mathrm{Bin}(2r,\frac12)$ variable conditioned on being even,
respectively odd. Their moments of order $\alpha$ differ by a multiple of
$\sum_{k=0}^{2r}(-1)^k\binom{2r}kk^\alpha$, which is zero for $\alpha<2r$ as a
finite difference of order $2r$ of a polynomial of degree $\alpha$. Thus the
moments of $s$ on the two sides agree through order $2r-1$.

All state atoms lie in $[0,rh]$, and all control atoms lie in $\{0,1\}$. With
cells as in the statement, $s$ carries all its mass, on both sides, in the
cell that contains $[0,rh]$, so the cellwise conditions at $\{s\}$ reduce to
the moment equalities just proved. Hence $\mu$ is feasible for the relaxation
of order $2r-1$. By the symmetry $k\mapsto2r-k$, $\E\bar u=r/2$ and
$\E\bar v=(r-1)/2$. Since $F_{r,h}=\lambda(\bar u+\bar v)$ on the supports,
the value of $\mu$ is $\lambda(r-\frac12)=\frac14-\frac{2r+1}{16r}$. Each state
variable occupies only the cell that contains $[0,rh]$, and each control only
the cells that contain $0$ and $1$; with $2r-1$ variables of each kind this
gives the bound on $\sum_i\omega_i(\mu)^2$.

\emph{Consequence.} Given $k$, choose $r$ with $2r-1\ge k$, so that $\mu$ is
also feasible for the relaxation of order $k$, and fix $p=3$, $\bar\nu=2$ and
$g=g_r$. The value of $\mu$ is $\OPT-\frac{2r+1}{16r}<\OPT-\frac18$, whereas
$\sum_i\omega_i(\mu)^2\to0$ as $h,h'\to0$ with $r$ fixed.
\end{proof}

The two bag chains realize the even and the odd part of a binomial law on the
separator, whose moments agree up to order $2r-1$. The example does not
contradict an error bound in terms of $L$, which is of order $r/h^2$ here. It
also does not survive exact matching of separator laws: on a tree
decomposition, bag laws that agree on every separator glue to a probability
measure on $X$ \cite{Vorobev1962}, whose value is at least $\OPT$. The
following remark shows that an inequality coupling both sides of a separator
detects the gap.

\begin{remark}[The smallest case]\label{lim:rem:moments}
The case $r=2$ admits a four-variable form. On $[0,1]^4$ with
$0<h\le\frac12$, let
\[
F_h=8\Bigl(\frac sh-\frac{u_1+u_2}2\Bigr)^2+8\Bigl(\frac sh-\frac14-\frac v2
\Bigr)^2+u_1(1-u_1)+u_2(1-u_2)+v(1-v)+\frac{u_1+u_2+v}{16},
\]
with bags $\{s,u_1,u_2\}$ and $\{s,v\}$. Direct expansion with
$\delta=s/h-\frac18-\frac{u_1+u_2+v}4$ gives
$F_h-\frac14=16\delta^2+2[(1-v)u_1u_2+v(1-u_1)(1-u_2)]+(u_1+u_2+v)/16$.
Hence the minimizer $(h/8,0,0,0)$ is unique and $g=1/22$ is valid, using
$\norm{x-x^*}^2\le2\delta^2+\frac{11}8\norm{(u_1,u_2,v)}^2$ and
$u_1+u_2+v\ge\norm{(u_1,u_2,v)}^2$. Moreover
$\nabla^2F_h+2I=16\chi_1\chi_1^{\T}+16\chi_2\chi_2^{\T}+2e_se_s^{\T}$ with
$\chi_1=(1/h,-\frac12,-\frac12,0)$ and $\chi_2=(1/h,0,0,-\frac12)$. Since
$(0,1,-1,0)$ is orthogonal to $\chi_1$, $\chi_2$ and $e_s$, $\nu=2$ exactly,
while $L=32/h^2$. The left measure with masses
$\frac18,\frac38,\frac38,\frac18$ at $(s/h,u_1,u_2)=(0,0,0),(\frac12,1,0),
(\frac12,0,1),(1,1,1)$ and the right measure with masses $\frac12,\frac12$
at $(s/h,v)=(\frac14,0),(\frac34,1)$ have equal moments of $s/h$ through order
three ($\frac12,\frac5{16},\frac7{32}$) but not four ($\frac{11}{64}$ versus
$\frac{41}{256}$). Their value is $3/32$, a gap of $5/32$.

These bag moments admit a global positive semidefinite completion.
With mean $(h/2,\frac12,\frac12,\frac12)$ and covariance
$\frac1{16}aa^{\T}+\frac3{16}bb^{\T}$, $a=(h,1,1,2)$,
$b=(0,1,-1,0)$, the completion satisfies every McCormick inequality for
$s\in[0,2h]$ and the control boxes. It violates the valid inequality
$u_1u_2+v-u_1v-u_2v=(1-v)u_1u_2+v(1-u_1)(1-u_2)\ge0$, a triangle inequality
of the Boolean quadric polytope \cite{Padberg1989}, with pseudo-expectation
$-1/8$. This inequality couples variables on both sides of the separator.
\end{remark}
